\chapter{Índice de notación}
Cada página remite al lugar donde se introduce o se utiliza sistemáticamente el símbolo. Las normas de funciones se toman respecto de la medida indicada en el contexto; en el círculo usamos medida de masa total uno. Las convenciones de signo de Fourier se muestran por separado.
\small
\begin{longtable}{@{}p{.23\textwidth}p{.64\textwidth}r@{}}
\toprule Símbolo & Significado & Pág.\\\midrule
\endfirsthead
\toprule Símbolo & Significado & Pág.\\\midrule
\endhead
\bottomrule\endfoot
$\N,\Z,\R,\C$ & Naturales positivos, enteros, reales y complejos & \pageref{cap:v4-01}\\
$s_n,\sigma_n$ & Suma parcial y promedio de Cesàro & \pageref{cap:v4-01}\\
$\T$ & Círculo unitario, también parametrizado por $t$ módulo $2\pi$ & \pageref{cap:v4-02}\\
$dm(t)$ & Medida normalizada $dt/(2\pi)$ sobre el círculo & \pageref{cap:v4-02}\\
$L^p,\norm f_p$ & Clases de funciones integrables a la potencia $p$ y su norma & \pageref{cap:v4-02}\\
$\ip uv$ & Producto interno, lineal en $u$ & \pageref{cap:v4-02}\\
$e_n(t)$ & Carácter periódico $e^{int}$ & \pageref{cap:v4-02}\\
$c_n(f)$ & Coeficiente $\int_\T f(t)e^{-int}dm(t)$ & \pageref{cap:v4-02}\\
$S_Nf$ & Suma parcial de Fourier, de $-N$ a $N$ & \pageref{cap:v4-03}\\
$D_N$ & Núcleo de Dirichlet $\sum_{|n|\le N}e^{int}$ & \pageref{cap:v4-03}\\
$K_N$ & Núcleo de Fejér $N^{-1}\sum_{j=0}^{N-1}D_j$ & \pageref{cap:v4-03}\\
$\sigma_N f$ & Promedio $N^{-1}\sum_{j=0}^{N-1}S_jf$ & \pageref{cap:v4-03}\\
$\Cc(X)$ & Funciones continuas de soporte compacto & \pageref{cap:v4-04}\\
$C_0(X)$ & Funciones continuas que se anulan en el infinito & \pageref{cap:v4-07}\\
$f*g$ & Convolución, con la medida de Haar especificada & \pageref{cap:v4-05}\\
$k_j,e_U$ & Núcleos de aproximación de la identidad & \pageref{cap:v4-05}\\
$\Z/q\Z$ & Grupo cíclico finito, con adición módulo $q$ & \pageref{cap:v4-05}\\
$\widehat a,\mathcal F_q$ & Transformada finita y su versión unitaria $q^{-1/2}\widehat a$ & \pageref{cap:v4-05}\\
$\widehat f(\xi)$ & Transformada euclidiana $\int_{\R^d}f(x)e^{ix\cdot\xi}dx$ & \pageref{cap:v4-07}\\
$\mathcal F f$ & Transformada unitaria $(2\pi)^{-d/2}\widehat f$ en $\R^d$ & \pageref{cap:v4-08}\\
$\Sscr(\R^d)$ & Espacio de Schwartz & \pageref{cap:v4-07}\\
$\alpha,D^\alpha$ & Multiíndice y derivada parcial correspondiente & \pageref{cap:v4-07}\\
$H_A$ & Subespacio de $L^2$ cuyo transformado se anula en $A$ & \pageref{cap:v4-08}\\
$PW_a$ & Funciones de $L^2(\R)$ con frecuencias en $[-a,a]$ & \pageref{cap:v4-08}\\
$K_y$ & Núcleo reproductor en el punto $y$ de $PW_a$ & \pageref{cap:v4-08}\\
$\operatorname{sinc}(t)$ & $\sin(\pi t)/(\pi t)$, extendida con valor $1$ en cero & \pageref{cap:v4-08}\\
$G/H,q$ & Espacio de clases izquierdas y proyección canónica & \pageref{cap:v4-09}\\
$G_0$ & Componente conexa de la identidad & \pageref{cap:v4-09}\\
$dx,\mu$ & Medida de Haar izquierda fijada & \pageref{cap:v4-10}\\
$(f:\phi)$ & Ínfimo de coeficientes de recubrimientos por traslaciones & \pageref{v4:haar-recubrimientos}\\
$\Delta_G,\Delta$ & Función modular: $\mu(Aa)=\Delta(a)\mu(A)$ & \pageref{v4:mod-convencion}\\
$L_af,R_af$ & Traslaciones $f(a^{-1}x)$ y $f(xa)$ & \pageref{cap:v4-11}\\
$T_Hf$ & Integral de $f$ a lo largo de las fibras de $G/H$ & \pageref{v4:mod-promedio}\\
$f^*$ & Involución $\Delta(x^{-1})\overline{f(x^{-1})}$ & \pageref{cap:v4-12}\\
$\mathcal U(H)$ & Grupo de operadores unitarios en $H$ & \pageref{cap:v4-13}\\
$\lambda,\rho$ & Representaciones regulares izquierda y derecha & \pageref{v4:rep-regulares}\\
$\Bcal(H,H')$ & Operadores lineales acotados de $H$ en $H'$ & \pageref{cap:v4-13}\\
$C(\pi,\pi'),C(\pi)$ & Entrelazadores y conmutante de una representación & \pageref{cap:v4-13}\\
$\Pi_U(f)$ & Representación integrada $\int_G f(x)U(x)dx$ & \pageref{v4-l24-p:integral}\\
$C^*(G),C_r^*(G)$ & Álgebra universal y álgebra reducida del grupo & \pageref{v4-l24-d:cstar}\\
$\widehat G$ & Grupo de caracteres continuos $G\to\T$ & \pageref{cap:v4-14}\\
$[K,U]$ & Subbásico de la topología compacto-abierta & \pageref{cap:v4-14}\\
$\widehat f(\chi)$ & Transformada abstracta $\int_Gf(x)\chi(x)dx$ & \pageref{cap:v4-15}\\
$\Delta(A)$ & Espacio de caracteres de un álgebra conmutativa & \pageref{cap:v4-15}\\
$\widehat\mu$ & Haar en $\widehat G$ normalizada por Plancherel & \pageref{v4:haar-dual}\\
$\widehat{\widehat G}$ & Bidual del grupo topológico $G$ & \pageref{v4:pontryagin}\\
$H^\perp$ & Aniquilador: caracteres de $G$ iguales a $1$ sobre $H$ & \pageref{v4:doble-aniquilador}\\
\end{longtable}
\normalsize
